\chapter{Bẫy Môi Trường Sai}
\markboth{Bẫy Môi Trường Sai}{The Trap of the Wrong Environment}

\section{Ở Nơi Chuyện Xấu Thành Bình Thường}
\begin{truyen}{Ở Nơi Chuyện Xấu Thành Bình Thường}{Living Where the Wrong Thing Feels Normal}
\chuhoa{C}{ó người trẻ bước vào một môi trường mà nói dối nhẹ, trễ hẹn, làm ẩu, hay xem thường người khác đều được coi là chuyện thường.}
\textit{(Some young people enter environments where small lies, lateness, careless work, and disrespect are treated as normal.)}

Ban đầu họ còn khó chịu.
\textit{(At first, they feel uncomfortable.)}

Nhưng sống lâu trong đó, ngưỡng khó chịu bị bào mòn dần và điều lệch bắt đầu bớt lệch trong mắt mình.
\textit{(But after living there for a while, their inner alarm wears down and the unhealthy thing begins to feel less wrong.)}

Môi trường có thể không ép em thành xấu ngay, nhưng nó dạy em quen với cái xấu nhanh hơn em tưởng.
\textit{(An environment may not force you to become unhealthy immediately, but it trains you to get used to it faster than you think.)}

\end{truyen}

\begin{baihoc}
Có những thứ em không chủ động chọn, nhưng ở cạnh quá lâu thì vẫn ngấm vào mình.
\textit{(Some things are not chosen on purpose, yet they still soak into you if you stay beside them too long.)}
\end{baihoc}

\begin{ghinhoanh}
\textbf{Short English to remember:}

Some things still soak into you if you stay beside them too long.

\end{ghinhoanh}

\begin{tuhocnhanh}
\textbf{Câu hỏi tự học / Self-study questions}

1. Điều gì đang bị bình thường hóa? \textit{What is being normalized?}

2. Ngưỡng nào đang mòn đi? \textit{What inner standard is wearing down?}

3. Em đang ở cạnh điều gì quá lâu? \textit{What have you been staying beside for too long?}
\end{tuhocnhanh}
\ngancach

\section{Chơi Với Người Luôn Kéo Mình Xuống}
\begin{truyen}{Chơi Với Người Luôn Kéo Mình Xuống}{Staying Close to People Who Keep Pulling You Down}
\chuhoa{C}{ó nhóm bạn mà hễ ai muốn học hơn, sống kỷ luật hơn, hay làm điều nghiêm túc hơn là lập tức bị chọc cho mềm đi.}
\textit{(Some circles immediately mock anyone who tries to study harder, live more disciplined, or take life more seriously.)}

Những câu đùa tưởng nhẹ nhưng lặp lại đủ nhiều sẽ làm ý chí xẹp xuống.
\textit{(Those jokes may seem light, but repeated often enough they flatten determination.)}

Không phải ai kéo em cười cũng đang kéo em lên.
\textit{(Not everyone who makes you laugh is pulling you upward.)}

Một môi trường luôn xem nỗ lực là lố bịch sẽ rất khó cho người trẻ giữ được đường dài.
\textit{(An environment that treats effort as ridiculous makes long-term growth very hard.)}

\end{truyen}

\begin{baihoc}
Người đi đường dài cần ở gần những người không chế nhạo phần nghiêm túc của mình.
\textit{(People who want a long road need to stay near those who do not mock their seriousness.)}
\end{baihoc}

\begin{ghinhoanh}
\textbf{Short English to remember:}

Stay near people who do not mock your seriousness.

\end{ghinhoanh}

\begin{tuhocnhanh}
\textbf{Câu hỏi tự học / Self-study questions}

1. Nhóm này đang làm yếu điều gì? \textit{What is this group weakening?}

2. Vì sao những câu đùa lại có sức nặng? \textit{Why can jokes carry such weight?}

3. Em cần ở gần kiểu người nào hơn? \textit{What kind of people do you need to stay closer to?}
\end{tuhocnhanh}
\ngancach

\section{Làm Ở Nơi Không Có Chuẩn Để Học}
\begin{truyen}{Làm Ở Nơi Không Có Chuẩn Để Học}{Working Where There Is No Standard to Learn From}
\chuhoa{C}{ó những nơi không ai thực sự giỏi hơn để mình học, không ai giữ chuẩn đủ cao để mình noi theo.}
\textit{(Some places offer no one truly better to learn from and no standard high enough to imitate.)}

Ngày qua ngày, một người trẻ vẫn làm việc, vẫn bận, vẫn thấy mình đang lớn.
\textit{(Day after day, a young person keeps working, stays busy, and feels as if they are growing.)}

Nhưng thiếu chuẩn tốt, họ rất dễ nhầm giữa quen việc và tiến bộ.
\textit{(But without good standards, it becomes easy to confuse familiarity with progress.)}

Môi trường yếu không chỉ cho ít cơ hội. Nó còn cho ảo giác rằng mức hiện tại đã là đủ.
\textit{(A weak environment does not only offer fewer opportunities. It also creates the illusion that the current level is enough.)}

\end{truyen}

\begin{baihoc}
Muốn lớn nhanh, em cần ít nhất một nơi hoặc một người giữ chuẩn cao hơn em hiện tại.
\textit{(To grow faster, you need at least one place or one person whose standard is higher than your current self.)}
\end{baihoc}

\begin{ghinhoanh}
\textbf{Short English to remember:}

To grow faster, you need a standard higher than your current self.

\end{ghinhoanh}

\begin{tuhocnhanh}
\textbf{Câu hỏi tự học / Self-study questions}

1. Môi trường này đang thiếu điều gì? \textit{What is this environment missing?}

2. Ảo giác nào dễ xuất hiện? \textit{What illusion appears easily here?}

3. Em đang học chuẩn từ ai hoặc từ đâu? \textit{Whose standard are you learning from?}
\end{tuhocnhanh}
\ngancach

\section{Biết Không Hợp Mà Vẫn Ở Vì Quen}
\begin{truyen}{Biết Không Hợp Mà Vẫn Ở Vì Quen}{Knowing It Does Not Fit but Staying Anyway}
\chuhoa{N}{hiều người sớm cảm thấy môi trường hiện tại không hợp giá trị, không hợp định hướng, không hợp kiểu lớn lên mình muốn.}
\textit{(Many people sense early that their environment does not fit their values, direction, or the kind of growth they want.)}

Nhưng vì quen người, quen lịch, quen cách sống, họ vẫn ở lại rất lâu.
\textit{(But because they are used to the people, the rhythm, and the lifestyle, they stay for a very long time.)}

Ở lại trong một nơi sai không phải lúc nào cũng tạo ra bi kịch ngay.
\textit{(Staying in the wrong place does not always create immediate tragedy.)}

Nó thường lấy đi dần dần: một ít nhiệt, một ít tin tưởng, một ít bản sắc.
\textit{(It usually takes things away gradually: some fire, some confidence, some identity.)}

\end{truyen}

\begin{baihoc}
Sai môi trường hiếm khi hủy em một lần. Nó mài em đi từng chút.
\textit{(The wrong environment rarely destroys you all at once. It wears you down a little at a time.)}
\end{baihoc}

\begin{ghinhoanh}
\textbf{Short English to remember:}

The wrong environment wears you down a little at a time.

\end{ghinhoanh}

\begin{tuhocnhanh}
\textbf{Câu hỏi tự học / Self-study questions}

1. Vì sao nhân vật vẫn ở lại? \textit{Why does the character stay?}

2. Điều gì đang bị lấy đi dần? \textit{What is being slowly taken away?}

3. Em có cần rời một môi trường nào không? \textit{Do you need to leave an environment?}
\end{tuhocnhanh}
\ngancach

\section{Cười Theo Điều Mình Biết Là Sai}
\begin{truyen}{Cười Theo Điều Mình Biết Là Sai}{Laughing Along with What You Know Is Wrong}
\chuhoa{C}{ó người ở trong một nhóm mà nhiều trò đùa, cách nói, hay thói quen đều lệch đi, nhưng vẫn cười theo để khỏi lạc lõng.}
\textit{(Some people stay in groups where many jokes, comments, or habits go wrong, yet they still laugh along just to avoid feeling left out.)}

Ban đầu họ nghĩ đó chỉ là chuyện nhỏ.
\textit{(At first, they think it is only a small matter.)}

Nhưng điều mình cười theo nhiều lần rất dễ thành điều mình dần cho phép tồn tại bình thường.
\textit{(But what you repeatedly laugh along with slowly becomes something you allow as normal.)}

Môi trường sai hiếm khi bẻ mình một cú thật mạnh. Nó thường mài nhẹ chuẩn của mình từng chút.
\textit{(A wrong environment rarely breaks you in one violent blow. It usually wears your standards down little by little.)}

\end{truyen}

\begin{baihoc}
Đừng coi nhẹ những điều em đang phải cười theo để được ở lại.
\textit{(Do not dismiss the things you have to laugh along with just to stay.)}
\end{baihoc}

\begin{ghinhoanh}
\textbf{Short English to remember:}

What you laugh along with can slowly reshape your standards.

\end{ghinhoanh}

\begin{tuhocnhanh}
\textbf{Câu hỏi tự học / Self-study questions}

1. Nhân vật đang cười theo điều gì? \textit{What is the character laughing along with?}

2. Điều gì đang bị mài mòn? \textit{What is being worn down?}

3. Em đang phải làm ngơ trước điều sai nào để hòa nhập? \textit{What wrong thing are you ignoring just to fit in?}
\end{tuhocnhanh}
\ngancach

\section{Ở Gần Người Coi Thường Nỗ Lực}
\begin{truyen}{Ở Gần Người Coi Thường Nỗ Lực}{Staying Close to People Who Mock Effort}
\chuhoa{C}{ó người ở trong một nhóm mà cứ ai học nghiêm túc, làm chăm, hay sống có kỷ luật là bị trêu là màu mè hoặc làm quá.}
\textit{(Some people stay in groups where anyone serious, diligent, or disciplined gets mocked as excessive or fake.)}

Ban đầu họ chỉ cười cho qua.
\textit{(At first, they laugh it off.)}

Nhưng lâu dần, chính họ cũng ngại cố gắng vì sợ bị nhìn lạc nhịp với nhóm.
\textit{(But over time, they themselves become afraid to try because they do not want to seem out of step.)}

Môi trường coi thường nỗ lực sẽ khiến phần tốt đẹp trong em phải giấu đi để tồn tại.
\textit{(An environment that mocks effort forces the best part of you to hide in order to survive.)}

\end{truyen}

\begin{baihoc}
Nơi nào làm em ngại cố gắng là nơi rất đáng cảnh giác.
\textit{(Any place that makes you ashamed of effort deserves caution.)}
\end{baihoc}

\begin{ghinhoanh}
\textbf{Short English to remember:}

Beware of any place that makes effort feel shameful.

\end{ghinhoanh}

\begin{tuhocnhanh}
\textbf{Câu hỏi tự học / Self-study questions}

1. Nhóm này đang phản ứng ra sao với nỗ lực? \textit{How does this group react to effort?}

2. Điều gì trong em phải giấu đi? \textit{What in you has to be hidden?}

3. Em có đang ngại sống nghiêm túc vì môi trường quanh mình không? \textit{Are you ashamed of being serious because of your environment?}
\end{tuhocnhanh}
\ngancach

\section{Nhóm Sống Cạn Kéo Mơ Ước Cạn Dần}
\begin{truyen}{Nhóm Sống Cạn Kéo Mơ Ước Cạn Dần}{A Shallow Circle Can Thin Out Your Dreams}
\chuhoa{C}{ó người lúc đầu vẫn có ước mơ và chuẩn sống riêng, nhưng ở lâu trong một nhóm chỉ quan tâm chuyện ngắn hạn thì dần cũng nguội đi.}
\textit{(Some people start with real dreams and standards, but after staying too long in a group that only cares about short-term things, they slowly cool down.)}

Họ vẫn cười nói bình thường, nhưng bên trong ít còn lửa như trước.
\textit{(They still smile and speak normally, but inside there is less fire than before.)}

Môi trường không chỉ ảnh hưởng hành vi. Nó còn ảnh hưởng biên độ của điều em dám mong về cuộc đời mình.
\textit{(Environment shapes not only behavior. It also shapes the range of what you dare to hope for in life.)}

Ở lâu nơi sống cạn, mình rất dễ quên rằng mình từng muốn sống sâu hơn thế.
\textit{(Stay too long in shallow company, and you may forget you once wanted a deeper life.)}

\end{truyen}

\begin{baihoc}
Nơi em ở lâu sẽ dạy em mơ tới mức nào.
\textit{(The place you stay in longest will teach you how far you are allowed to dream.)}
\end{baihoc}

\begin{ghinhoanh}
\textbf{Short English to remember:}

Your environment quietly teaches you how far to dream.

\end{ghinhoanh}

\begin{tuhocnhanh}
\textbf{Câu hỏi tự học / Self-study questions}

1. Điều gì trong nhân vật đang nguội đi? \textit{What in the character is cooling down?}

2. Môi trường này đang dạy mức mơ nào? \textit{What level of dreaming is this environment teaching?}

3. Em có cần ở gần những người sống sâu hơn không? \textit{Do you need to stay closer to people who live more deeply?}
\end{tuhocnhanh}
\ngancach

\section{Môi Trường Thưởng Cho Sự Cẩu Thả}
\begin{truyen}{Môi Trường Thưởng Cho Sự Cẩu Thả}{An Environment That Rewards Carelessness}
\chuhoa{C}{ó nơi làm qua loa vẫn được chấp nhận, hứa không giữ cũng chẳng sao, chất lượng thấp cũng không ai nhắc.}
\textit{(Some places accept sloppy work, broken promises, and low quality without consequence.)}

Ở một môi trường như vậy, người nghiêm túc rất dễ mỏi hoặc bị xem là khó tính.
\textit{(In such an environment, serious people quickly grow tired or are seen as difficult.)}

Nếu ở lâu, mình hoặc sẽ phải chiến đấu liên tục, hoặc sẽ bị lây thói cẩu thả lúc nào không hay.
\textit{(If you stay too long, you either fight constantly or slowly catch the habit of carelessness.)}

Chuẩn sống của con người không chỉ cần ý chí. Nó còn cần một bầu khí đủ lành để thở.
\textit{(A person's standards need more than willpower. They also need an atmosphere healthy enough to breathe in.)}

\end{truyen}

\begin{baihoc}
Sống đúng trong chỗ sai quá lâu sẽ làm em kiệt hoặc làm em lệch.
\textit{(Staying right in a wrong place too long will either exhaust you or bend you.)}
\end{baihoc}

\begin{ghinhoanh}
\textbf{Short English to remember:}

Stay too long in a wrong place, and you either exhaust or bend.

\end{ghinhoanh}

\begin{tuhocnhanh}
\textbf{Câu hỏi tự học / Self-study questions}

1. Nơi này đang thưởng cho điều gì? \textit{What is this place rewarding?}

2. Nếu ở lâu sẽ có hai nguy cơ nào? \textit{What two risks appear if you stay too long?}

3. Chuẩn nào của em đang cần một môi trường lành hơn? \textit{What standard in your life needs a healthier environment?}
\end{tuhocnhanh}
\ngancach

\section{Sợ Cô Đơn Nên Không Dám Rời Chỗ Sai}
\begin{truyen}{Sợ Cô Đơn Nên Không Dám Rời Chỗ Sai}{Fearing Loneliness and Staying in the Wrong Place}
\chuhoa{C}{ó người biết rõ chỗ mình đang ở không còn hợp, nhưng vẫn nấn ná vì sợ khoảng trống sau khi rời ra.}
\textit{(Some people know clearly that where they are no longer fits, yet they linger because they fear the emptiness that comes after leaving.)}

Họ chọn một sự quen sai còn hơn một khoảng lặng chưa biết sẽ ra sao.
\textit{(They choose familiar wrongness over an unknown silence.)}

Nhưng có những khoảng cô đơn lành hơn rất nhiều một sự thuộc về làm mình nhỏ dần.
\textit{(But some forms of loneliness are far healthier than belonging somewhere that slowly shrinks you.)}

Không dám rời chỗ sai vì sợ cô đơn thường chỉ đổi một nỗi sợ ngắn lấy một sự hao mòn dài.
\textit{(Refusing to leave a wrong place out of fear of loneliness often trades short fear for long erosion.)}

\end{truyen}

\begin{baihoc}
Cô đơn tạm thời có khi là cánh cửa đi vào một môi trường đúng hơn.
\textit{(Temporary loneliness can be the doorway to a healthier environment.)}
\end{baihoc}

\begin{ghinhoanh}
\textbf{Short English to remember:}

Temporary loneliness may open the door to a better place.

\end{ghinhoanh}

\begin{tuhocnhanh}
\textbf{Câu hỏi tự học / Self-study questions}

1. Điều gì giữ nhân vật ở lại? \textit{What keeps the character staying?}

2. Họ đang đổi điều gì lấy điều gì? \textit{What are they trading for what?}

3. Em có đang sợ một khoảng trống cần thiết không? \textit{Are you afraid of a necessary empty space?}
\end{tuhocnhanh}
\ngancach

\section{Chọn Môi Trường Nâng Chuẩn Mình}
\begin{truyen}{Chọn Môi Trường Nâng Chuẩn Mình}{Choosing an Environment That Raises Your Standards}
\chuhoa{C}{ó người đổi đời không phải vì một bí quyết lạ, mà vì dũng cảm đổi chỗ đứng sang nơi mọi người sống nghiêm hơn, tử tế hơn, và có chuẩn hơn.}
\textit{(Some people change their life not through a secret formula, but by bravely moving into a place where others live with more discipline, kindness, and standards.)}

Ban đầu họ thấy mình nhỏ và hơi lạc lõng.
\textit{(At first, they feel small and slightly out of place.)}

Nhưng chính cảm giác đó lại kéo họ đi lên từng chút, vì bầu khí mới nâng đỡ phần tốt trong họ.
\textit{(But that very feeling pulls them upward little by little because the new atmosphere supports their better part.)}

Một môi trường đúng không làm em dễ chịu ngay lập tức. Nó làm em tốt hơn dần dần.
\textit{(A right environment does not always make you comfortable immediately. It makes you better gradually.)}

\end{truyen}

\begin{baihoc}
Chọn đúng môi trường là một cách rút ngắn nhiều vòng tự đấu với chính mình.
\textit{(Choosing the right environment shortens many rounds of fighting yourself.)}
\end{baihoc}

\begin{ghinhoanh}
\textbf{Short English to remember:}

The right environment shortens many inner battles.

\end{ghinhoanh}

\begin{tuhocnhanh}
\textbf{Câu hỏi tự học / Self-study questions}

1. Điều gì thay đổi khi nhân vật đổi môi trường? \textit{What changes when the character changes environments?}

2. Vì sao ban đầu lại thấy lạc lõng? \textit{Why does it feel uncomfortable at first?}

3. Em cần tiến gần hơn tới môi trường nào? \textit{What kind of environment do you need to move closer to?}
\end{tuhocnhanh}
